% Einheit 3 von 5 – Brüche vergleichen (bank/brueche-dezimalzahlen/e3.jsonl, zusammenbau v0.1)
%% TODO Zweigzeile Teil 1: was hier gelernt wird (Satz in Schülersprache aus dem Einheitstitel) – Einheit 3
\einheitenkopf[e3]{Einheit 3 von 5 · Brüche vergleichen}
\zweigzeile{ab Kl. 5, je nach Buch bis Kl. 6 · P10 oft}
\setcounter{aufgabe}{17}

%% TODO Titel: Ich-kann-Satz fehlt in der Bank (Platzhalter: Kettenname) – Nr. 18
%% TODO Anweisung: ein Satz über der ersten Teilaufgabe fehlt in der Bank – Nr. 18
\begin{aufgabe}{Vergleichen}
%% TODO \rechenplatz{n}: Schrittzahl des Grundfalls fehlt in der Bank (nur bei fester Schreibform, 2.3 b)
\begin{teile}
\teil $\frac{3}{8}$ und $\frac{5}{8}$ – welcher Vergleichsweg ist am schnellsten? Kreuze an.\\ \kreuz{gleicher Nenner}\\ \kreuz{gleicher Zähler}\\ \kreuz{mit einem Halben vergleichen}\\ \kreuz{gleichnamig machen}
\teil Welcher Bruch ist größer: $\frac{2}{11}$ oder $\frac{7}{11}$? \leerfeld
\teil Welcher Bruch ist größer: $\frac{3}{8}$ oder $\frac{3}{5}$? \leerfeld
\teil Welcher Bruch ist größer: $\frac{4}{9}$ oder $\frac{5}{8}$? Vergleiche mit $\frac{1}{2}$. \leerfeld
\teil Welcher Bruch ist größer: $\frac{3}{4}$ oder $\frac{5}{7}$? \leerfeld
\teil Trage $\frac{3}{8}$ und $\frac{7}{8}$ am Zahlenstrahl ein.

\zahlenstrahl[xmin=0,xmax=1,xstep=0.125,karo=1.2]{0/0, 1/1}
\teil Ordne von klein nach groß: $\frac{2}{3}$, $\frac{1}{6}$, $\frac{5}{6}$. \leerfeld $<$ \leerfeld $<$ \leerfeld
\teil Gib eine Zahl an, die zwischen $\frac{1}{4}$ und $\frac{3}{5}$ liegt. \hfill (P10 2014 OS) \\ \leerfeld
\end{teile}
\end{aufgabe}

%% TODO Titel: Ich-kann-Satz fehlt in der Bank (Platzhalter: Kettenname) – Nr. 19
%% TODO Anweisung: ein Satz über der ersten Teilaufgabe fehlt in der Bank – Nr. 19
\begin{aufgabe}{Zahl zwischen zwei Brüchen angeben}
\begin{teile}
\teil Gib einen Bruch zwischen $\frac{1}{5}$ und $\frac{2}{5}$ an. \leerfeld
\end{teile}
\end{aufgabe}

%% TODO Titel: Ich-kann-Satz fehlt in der Bank (Platzhalter: Kettenname) – Nr. 20
%% TODO Anweisung: ein Satz über der ersten Teilaufgabe fehlt in der Bank – Nr. 20
\begin{aufgabe}{Vergleichen – Fehler finden, Begründen, Darstellungswechsel, Anwendung}
\begin{teile}
\teil Kai sagt: $\frac{2}{3}$ und $\frac{5}{6}$ sind gleich groß, denn bei beiden fehlt nur ein Stück zum Ganzen. Finde den Fehler.
\teil Warum ist $\frac{1}{5}$ größer als $\frac{1}{6}$? Begründe ohne Rechnung.
\teil Welcher Bruch gehört zu A?

\zahlenstrahl[xmin=0,xmax=1,xstep=0.25,karo=3]{0/0, 1/1, 0.75/A}
\leerfeld
\teil Lara und Emil laden gleich große Dateien. Bei Lara sind $\frac{5}{8}$ geladen, bei Emil $\frac{2}{3}$. Wer ist schon weiter?
\end{teile}
\end{aufgabe}
